\chapter{What calibrated exposure preserves at the population level}
\label{ch:dg-recovery}

Finite exposure changes a diffusion observation before measurement noise is
added. In particular, a boxcar average of a continuous signal is continuously
differentiable and has zero quadratic variation. The local geometry theorem
therefore cannot be applied directly to an exposed trace as if it were the
original diffusion coordinate. The question here is different: does the
complete exposed observation law still determine the ideal law from which
that trace came?

Under the assumptions below, the answer is yes, even after a known independent
additive noise channel. The recovery uses dense temporal population access
and an infinite limiting reconstruction. It does not provide a stable
finite-data inverse or recover local quadratic variation by differentiating
noisy samples. The constituent deconvolution and averaging arguments are
established mathematical tools; historical priority is not asserted.

\section{Detector and access model}

Let $Y=(Y_t)_{t\in\mathbb R}$ be a real stationary continuous process,
bounded by a finite deterministic constant, and write $\mu=\mathbb E Y_0$.
Assume
\begin{equation}
 \operatorname{Cov}(Y_0,Y_t)\longrightarrow0
 \quad\hbox{as }|t|\longrightarrow\infty.
 \label{eq:dg-covariance-mixing}
\end{equation}
This condition is satisfied by a smooth bounded coordinate of a stationary
compact connected uniformly elliptic diffusion. A known exposure width
$\tau>0$ produces the causal boxcar signal
\begin{equation}
 Z_t=(A_\tau Y)_t
      :=\frac1\tau\int_{t-\tau}^tY_s\,ds.
 \label{eq:dg-exposure}
\end{equation}
Let $\mathcal T$ be a known countable dense subset of the real time axis.
At each distinct $t\in\mathcal T$, the retained measurement is
\begin{equation}
 B_t=Z_t+\xi_t,
 \label{eq:dg-readout}
\end{equation}
where the $\xi_t$ are independent of $Y$, independent across distinct
timestamps, and have a known probability law $\nu$. The known noise law
need not have finite variance or a nowhere-zero characteristic function for
the population uniqueness statement. Its independence and calibration are
part of the detector contract.

The accessible research object in this section is the complete joint law of
$(B_t)_{t\in\mathcal T}$, rather than one finite record. Repeated samples
of a finite collection of timestamps do not automatically supply this
infinite family of joint distributions. The clock, exposure rule, and noise
law are common to every model being compared.

\section{Noise removal as an identifiability statement}

\begin{lemma}[Known additive convolution on a bounded latent law]
\label{lem:dg-bounded-convolution}
Let $P$ and $Q$ be compactly supported probability laws on $\mathbb R^k$,
and let $N$ be any probability law on $\mathbb R^k$. If
$P*N=Q*N$, then $P=Q$.
\end{lemma}

\begin{proof}
The characteristic function $\widehat N$ is continuous and equals one at
the origin, so it is nonzero on some neighborhood of zero. The convolution
identity implies $\widehat P=\widehat Q$ on that neighborhood. Compact
support makes both characteristic functions real analytic on
$\mathbb R^k$ (indeed entire after complexification). Analytic uniqueness
therefore extends the equality to all real arguments. Uniqueness of
characteristic functions gives $P=Q$.
\end{proof}

Apply this lemma at any finite set of distinct timestamps, with the known
noise law $N=\nu^{\otimes k}$. Equality of the noisy record laws implies
equality of the corresponding finite-dimensional laws of $Z$. Since $Z$ is
continuous, its values on the countable dense set determine its law on
$C(\mathbb R)$, equipped with uniform convergence on compact intervals.
Thus the complete noisy law identifies the complete exposure law.

This proof does not divide by the noise characteristic function at its
zeros. It uses bounded latent support and uniqueness of analytic
continuation; it supplies no stability estimate for that continuation. If
bounded support is dropped, a standard sufficient replacement is a noise
characteristic function that is nowhere zero. These are different
identifiability contracts. Known correlated additive noise can also be
treated if its complete joint laws are supplied and remain independent of
the signal; knowing only its one-time marginal would not suffice.

\subsection{What infinitely dense independent readouts can recover}

For the iid noise contract, dense access permits a stronger statement that
does not require bounded signal amplitude. Let $Z$ have continuous finite
real-valued paths and let $B_t=Z_t+\xi_t$ be observed on the countable dense
set, with the same known iid noise independent of $Z$. Then the complete
continuous exposure path is a measurable function of the infinitely dense
noisy record, almost surely.

To prove this, fix a target $t$ in the dense set and a deterministic sequence
of distinct available timestamps $t_j\to t$. For any fixed real $u$, the
bounded complex observations satisfy
\begin{equation}
 \frac1n\sum_{j=1}^n e^{iuB_{t_j}}
 \longrightarrow e^{iuZ_t}\widehat\nu(u)
 \quad\hbox{almost surely}.
 \label{eq:dg-dense-noise-average}
\end{equation}
Conditioned on the entire continuous $Z$ path, the centered summands are
independent bounded variables. The independent-variable strong law applies
because the sum of their variances divided by $j^2$ is finite. Their average
conditional means converge to the right-hand side by continuity of $Z$.
The argument uses no moment of $\xi_t$.

Choose a countable sequence $u_k\to0$ with $u_k\ne0$. The known noise
characteristic function is nonzero for all sufficiently large $k$.
Dividing the limits in~\eqref{eq:dg-dense-noise-average} by
$\widehat\nu(u_k)$ determines $e^{iu_kZ_t}$, and
\[
 Z_t=\lim_{k\to\infty}
       \frac{\operatorname{Im}(e^{iu_kZ_t})}{u_k}
\]
removes the phase ambiguity. A countable intersection of probability-one
events recovers every dense-set $Z_t$; continuity supplies the full path.
All operations are limits of measurable coordinate functions, with an
arbitrary default on the null set where the limits fail.

Consequently the main injectivity theorem below also holds with boundedness
replaced by finite second moment, using this dense iid reconstruction and
the same covariance-decay assumption. The compact diffusion application
remains bounded. The proposition spends infinitely many independent
readouts in arbitrarily small time neighborhoods. It is not a claim that
faster sampling makes actual sensor errors independent: overlapping
exposures, photon reuse, memory, or correlated readout noise can invalidate
the stipulated iid contract. The earlier bounded-law convolution argument
is still useful for known correlated noise at the population-law level,
where pathwise separation from one continuous record need not hold.

\section{An explicit finite-window reconstruction and its limit}

Continuity of $Y$ gives the pathwise identity
\begin{equation}
 \tau Z'_t=Y_t-Y_{t-\tau}.
 \label{eq:dg-derivative}
\end{equation}
The derivative is a measurable functional of the ideal $C^1$ exposure path,
for example the limit of rational difference quotients. For $N\geq1$ set
\begin{equation}
 \mathcal R_N Z(t)
 :=\mathbb E Z_0+
 \tau\sum_{k=0}^{N-1}\left(1-\frac{k}{N}\right)Z'_{t-k\tau}.
 \label{eq:dg-triangular-reconstruction}
\end{equation}
Telescoping gives the exact relation
\begin{equation}
 \mathcal R_N Z(t)
 =Y_t-\left\{\frac1N\sum_{j=1}^NY_{t-j\tau}-\mu\right\}.
 \label{eq:dg-reconstruction-error}
\end{equation}
Only a window whose length grows as $N\tau$ is used at each finite $N$,
but exact derivatives of the exposed path are required.

\begin{theorem}[Dense-law injectivity through boxcar exposure and known noise]
\label{thm:dg-injectivity}
Within the bounded continuous stationary class satisfying
Equation~\eqref{eq:dg-covariance-mixing}, equality of the noisy dense-time
laws in~\eqref{eq:dg-readout}, for a common known boxcar width and noise
law, implies equality of the ideal continuous path laws of $Y$.
\end{theorem}

\begin{proof}
The preceding lemma identifies the common continuous law of $Z$.
Stationarity gives $\mathbb E Z_0=\mu$. By
Equation~\eqref{eq:dg-reconstruction-error},
\begin{align}
 \mathbb E|\mathcal R_N Z(t)-Y_t|^2
 &=\frac1{N^2}\sum_{i,j=1}^N
          \operatorname{Cov}(Y_0,Y_{(i-j)\tau})\notag\\
 &=\frac{\operatorname{Var}(Y_0)}N
   +\frac2{N^2}\sum_{j=1}^{N-1}(N-j)
                \operatorname{Cov}(Y_0,Y_{j\tau})
 \longrightarrow0.
 \label{eq:dg-mean-square-recovery}
\end{align}
For the last limit, bound the finite initial part of the covariance sum by
a constant times $N^{-1}$, and use arbitrarily small covariance on its
remaining tail. The same conclusion holds jointly at every finite set of
times, by summing their mean-square errors. The law of each
$\mathcal R_N Z$ is fixed by the common exposure law, so the limiting
finite-dimensional laws of $Y$ agree. Continuity again identifies the full
path laws.
\end{proof}

An explicitly calibrated covariance-decay bound gives a finite-window
population estimate. If
$|\operatorname{Cov}(Y_0,Y_t)|\leq C e^{-\kappa|t|}$ for $C<\infty$
and $\kappa>0$, then
\begin{equation}
 \mathbb E|\mathcal R_N Z(t)-Y_t|^2
 \leq\frac{C}{N}\coth(\kappa\tau/2).
 \label{eq:dg-reconstruction-rate}
\end{equation}
This controls truncation of an ideal derivative reconstruction. It does
not control estimation of $Z'$ from finitely many noisy or exposed readings.
For example, if a derivative approximation has root mean-square error at
most $\delta$ at each of the retained times, and the stationary mean has
error at most $\delta_\mu$, the triangle inequality gives only
\[
 \|\widehat{\mathcal R_N Z}(t)-Y_t\|_2
 \leq \sqrt{\frac C N\coth(\kappa\tau/2)}
       +\frac{\tau(N+1)}2\,\delta+\delta_\mu.
\]
Increasing the reconstruction window can amplify derivative error. Any
finite-data use must pay for both terms under an explicit detector model.

\section{The spectral obstruction and a derivative-free population inverse}

The exposure multiplier at frequency $\lambda$ is
\begin{equation}
 H_\tau(\lambda)
 =e^{-i\lambda\tau/2}\frac{\sin(\lambda\tau/2)}{\lambda\tau/2},
 \qquad H_\tau(0)=1.
 \label{eq:dg-transfer}
\end{equation}
It vanishes exactly at $\lambda=2\pi k/\tau$, $k\in\mathbb Z\setminus\{0\}$.
The decay assumption excludes spectral atoms at those frequencies. The
minimal spectral premise for recovering a particular $Y$ is that its
spectral measure assign zero mass to those zeros; full mixing is a convenient
sufficient assumption.

For completeness, the usual Landweber reconstruction gives a second proof
that does not explicitly differentiate $Z$. On the stationary $L^2$ shift
space let $A=A_\tau$ and $T=A^*A$, so $0\preceq T\preceq I$. Define
\begin{equation}
 \mathcal L_N Z
   =\sum_{j=0}^{N-1}(I-T)^jA^*Z.
 \label{eq:dg-landweber}
\end{equation}
Since $Z=AY$,
\[
 \mathcal L_N Z=[I-(I-T)^N]Y.
\]
The spectral theorem gives $L^2$ convergence to $Y$ exactly when $Y$ has no
component in $\ker A$. The kernel consists of the nonzero boxcar-resonant
spectral modes above; the stationary mean is preserved because $H_\tau(0)=1$.
Each iterate is a finite linear combination of finite-window averages of
$Z$ and is therefore a fixed measurable function of its continuous path.

This inverse is also not uniformly stable as $N\to\infty$. Its operator
norm is at most $\sqrt N$, because for $0\leq x\leq1$,
\[
 \sqrt{x}\sum_{j=0}^{N-1}(1-x)^j
 =\frac{1-(1-x)^N}{\sqrt x}
 \leq\min\{N\sqrt x,x^{-1/2}\}\leq\sqrt N.
\]
At $x=0$ the expression is defined by its zero limit for fixed $N$.
For a stationary $L^2$ perturbation of the entire input process, this bounds
its amplification by $\sqrt N$. It is not an assertion that a finite
collection of independent point-noise values forms a directly integrable
continuous error process.

The spectral bias is explicitly
\[
 \|\mathcal L_N AY-Y\|_2^2
 =\int(1-|H_\tau(\lambda)|^2)^{2N}\,d\mu_Y(\lambda).
\]
Vanishing mass at exact zeros ensures convergence; a useful finite-error
rate additionally needs quantitative control of spectral mass near the
zeros and at large frequencies. Injectivity alone supplies no such rate.

\section{Consequences for the compatible diffusion minimum}

Every admissible compact connected uniformly elliptic hidden diffusion in
the observed-geometry chapter has a bounded smooth detector and exponentially
decaying centered correlations. Consequently, if its ideal coordinate is
observed through the detector contract here, the complete noisy dense-time
law determines its ideal continuous observation law.

In particular, fix the noisy exposed law generated by the stationary
cosine-circle process with parameters $\gamma>0$ and $\omega$, with the
same calibrated detector channel used in all compatible models. Every
compatible admissible hidden diffusion then has the same ideal law as that
circle. The earlier attained minimum therefore remains
\begin{equation}
 \min\sigma=\omega^2/\gamma.
 \label{eq:dg-minimum-preserved}
\end{equation}
The circle with the stated detector is feasible, so this is an attained
minimum, not only a lower bound. The result concerns entropy production in
the target hidden diffusion, not the unmodeled energy consumed by the
detector or the readout process.

The identification uses infinitely rich population information. It does
not show that a finite noisy experiment can attain the same accuracy as a
noiseless high-bandwidth experiment. In particular, $Z$ itself has zero
quadratic variation. Local geometry is recovered only after an appropriate
limiting reconstruction of the ideal law, rather than by applying the
previous local formula directly to $Z$ or $B$.

\section{When the width itself is identifiable}

Here and in the preceding compatible-minimum statement, the diffusion
assumptions are those of the earlier observed-geometry class: a smooth compact
connected manifold without boundary, smooth coefficients and strictly positive
stationary density, uniform ellipticity, and ordinary reversal of the complete
hidden state. No unspecified boundary condition or generalized reversal is
included by the shorter phrase ``compact elliptic'' below.

Within the compact elliptic model class, a known boxcar shape has a further
population property: its width can be identified from the exposed law even
if that width was not initially supplied. Assume the detector is nonconstant.
Let $f$ now denote the centered detector, and use the complex
$L^2(\pi)$ inner product. The symmetric part of the generator has a
Poincar\'e gap, so its mean-zero semigroup decays exponentially and the
resolvent exists at every $i\lambda$, $\lambda\in\mathbb R$.
For $u=(i\lambda-L)^{-1}f$, the ideal power spectral density is
\begin{equation}
 S_Y(\lambda)
 =2\operatorname{Re}\langle f,(i\lambda-L)^{-1}f\rangle_\pi
 =2\int_M\nabla\overline u^{\mathsf T}D\nabla u\,d\pi
 >0.
 \label{eq:dg-positive-spectrum}
\end{equation}
The identity follows from $f=(i\lambda-L)u$ and skew-adjointness of the
stationary current term. The Dirichlet form vanishes only for constant $u$;
in the mean-zero subspace this would force $u=0$ and then $f=0$, contrary
to the nonconstant detector assumption.

The centered exposed spectrum is
$S_Z(\lambda)=|H_\tau(\lambda)|^2S_Y(\lambda)$. Its zero set is
therefore exactly the nonzero boxcar frequencies. The smallest positive
zero equals $2\pi/\tau$ and identifies $\tau$. Known-noise deconvolution
may first be used to identify the exposure law; no finite-variance noise
assumption is needed if the bounded-law uniqueness argument is used.
This proposition assumes the boxcar shape, nonconstant diffusion detector,
continuous population spectrum and calibrated time. A constant signal has
no such width information. Estimating an exact spectral zero from a finite
noisy record is a separate problem, and an unknown general response kernel
is not identified by this argument.

\section{Necessary distinctions}

The mathematical ingredients have substantial prior art. Landweber's 1951
iteration is an established inverse method. Known-error deconvolution with
characteristic-function zeros is treated, under quantitative assumptions,
by Hall and Meister (2007). Nakakita and Uchida (2020) already estimate an
unknown exposure-to-frame ratio for convolutionally observed diffusions in
a shrinking-exposure, increasing-duration regime. The fixed-width
full-law argument here is a different observation regime, not the first
inference of an unknown shutter parameter. Markov-generator resolvent
formulas for power spectra are also established; Gupta and Khammash (2022)
provide a close explicit use for stochastic reaction networks. The focused
primary-source audit is recorded separately in this research phase. These
comparisons leave historical priority for any narrower combination open.

A stationary periodic signal with uniform random phase at frequency
$2\pi/\tau$ is erased exactly by the boxcar. Such a process is stationary
and can be ergodic under the full continuous time-shift group while failing
mixing and ergodicity under the sampled $\tau$-shift. Continuous-flow
ergodicity alone is therefore not a sufficient replacement for the
assumptions used here.

A fixed discrete sampling grid also supplies less information than the
dense-time law in the theorem. Continuity does not reconstruct an exposed
path between grid points. Exact aliasing or different hidden dynamics can
survive fixed-grid exposure data, and no claim of injectivity for that
access model follows from this chapter.

Finally, the inverse has no uniform stability bound over all admissible
mixing signals. Boxcar multipliers become arbitrarily small near their
zeros and at high frequency. Additional bandwidth, dynamical, detector or
resource restrictions are needed for a useful finite-error conclusion.
Population identifiability, physical attainability, finite-data confidence,
and experimental energy are separate claims throughout this argument.
